\begin{figure}[htbp]
	\centering
	\begin{tikzpicture}[>=Stealth, line width=0.8pt,
		scale=1.1]
		
		% ── Trave AB con carichi distribuiti ─────────────
		\draw[thick] (0,0) -- (9,0);
		\draw[dashed, thick] (0,-0.18) -- (9,-0.18);
		\fill (0,0) circle (2pt);
		\node[below=6pt] at (0,0) {$A$};
		\fill (9,0) circle (2pt);
		\node[below=6pt] at (9,0) {$B$};
		
		% p(z)
		\draw[thick, domain=0:9, samples=100]
		plot (\x, {1.2 + 0.5*sin(\x*40)});
		\foreach \x in {0.3, 0.9, 1.5, 2.1, 2.7, 3.3,
			3.9, 4.5, 5.1, 5.7, 6.3, 6.9,
			7.5, 8.1, 8.7}{
			\pgfmathsetmacro{\h}{1.2 + 0.5*sin(\x*40)}
			\draw[->, thin] (\x, \h) -- (\x, 0);
		}
		\node[right] at (9.1, {1.2+0.5*sin(360)}) {$p(z)$};
		
		% q(z)
		\foreach \x in {0.3, 0.9, 1.5, 2.1, 2.7, 3.3,
			3.9, 4.5, 5.1, 5.7, 6.3, 6.9,
			7.5, 8.1, 8.7}{
			\draw[->, thin] (\x-0.35, 0.35) -- (\x, 0.35);
		}
		\draw[thick] (0,0.35) -- (9,0.35);
		\node[right] at (9.1, 0.35) {$q(z)$};
		
		% c(z) — archi antiorari
		\foreach \x in {0.5, 1.3, 2.1, 2.9, 3.7, 4.5,
			5.3, 6.1, 6.9, 7.7, 8.5}{
			\draw[->] (\x+0.18, 0.72)
			arc[start angle=0, end angle=180,
			radius=0.18];
		}
		\node[right] at (9.1, 0.72) {$c(z)$};
		
		% Punti P e Q
		\draw[thin] (4.0, -0.3) -- (4.0, 0.1);
		\draw[thin] (5.2, -0.3) -- (5.2, 0.1);
		%\node[below=2pt] at (3.8, -0.3) {$P$};
		\node[below=2pt] at (4.0, -0.12) {\footnotesize$(z)$};
		%\node[below=2pt] at (5.4, -0.3) {$Q$};
		\node[below=2pt] at (5.3, -0.12)
		{\footnotesize$(z{+}\Delta z)$};
		\draw[|<->|, thin] (4.0,-0.75) -- (5.2,-0.75)
		node[midway, below=2pt] {$\Delta z$};
		
		% Linee tratteggiate verso concio
		\draw[dashed, gray!50, thin] (4.0,-0.18) -- (2.0,-2.0);
		\draw[dashed, gray!50, thin] (5.2,-0.18) -- (7.0,-2.0);
		
		% ── Concio ingrandito (più corto) ─────────────────
		\draw[thick] (2.0,-3.8) rectangle (7.0,-2.0);
		\draw[dashed, thick] (2.0,-4.0) -- (7.0,-4.0);
		
		% Centro del concio
		\coordinate (C) at (4.5,-2.9);
		
				
		% c dz: coppia antioraria sopra il centro
		\draw[->] (4.85,-2.8)
		arc[start angle=0, end angle=180,
		radius=0.38];
		\node[above=2pt, font=\small] at (4.5,-2.6) {$c(z)\,\Delta z$};
		
		% q dz: freccia da centro verso destra
		\draw[->, line width=1.2pt]
		(C) -- ++(1.5, 0)
		node[above=2pt, font=\small] {$q(z)\,\Delta z$};
		
		%p dz: freccia dal centro verso il basso
		\draw[->, line width=1.2pt]
		(C) -- ++(0, -0.9)
		node[pos=0.5, right=2pt, font=\small] {$p(z)\,\Delta z$};
		
		% ── Faccia sinistra (z) ─────────────────────────────
		% N
		\draw[->, line width=1.6pt]
		(2.0,-2.9) -- (0.7,-2.9);
		\node[above=2pt, font=\small] at (1.2,-2.9) {$N$};
		
		% T
		\draw[->, line width=1.6pt]
		(2.0,-2.9) -- (2.0,-1.7);
		\node[left=2pt, font=\small] at (2.0,-2.3) {$T$};
		
		% M
		\draw[->, line width=1.6pt]
		(0.7,-2.9) arc[start angle=180,
		end angle=90, radius=1.2];
		\node[left=2pt, font=\small] at (0.95,-1.9) {$M$};
		
		% ── Faccia destra (z+dz) ───────────────────────────────
		% N+dN
		\draw[->, line width=1.6pt]
		(7.0,-2.9) -- (8.3,-2.9);
		\node[above=2pt, font=\small] at (7.8,-2.9) {$N{+}\Delta N$};
		
		% T+dT
		\draw[->, line width=1.6pt]
		(7.0,-2.9) -- (7.0,-4.1);
		\node[right=2pt, font=\small] at (6.85,-3.4) {$T{+}\Delta T$};
		
		% M+dM
		\draw[->, line width=1.6pt]
		(7.0,-4.1) arc[start angle=270,
		end angle=360, radius=1.2];
		\node[right=2pt, font=\small] at (7.9,-3.6) {$M{+}\Delta M$};
		
	\end{tikzpicture}
	\caption{Concio elementare compreso fra le sezioni di
		ascissa $z$ e $z + \Delta z$ in un tratto regolare.
		Sulla faccia di ascissa $z$ le azioni interne agiscono
		con verso opposto a quello positivo, sulla faccia di
		ascissa $z + \Delta z$ con verso positivo e incrementate
		di $\Delta N$, $\Delta T$, $\Delta M$. Al primo ordine
		in $\Delta z$ le azioni distribuite si riducono alle
		risultanti $p(z)\,\Delta z$, $q(z)\,\Delta z$, $c(z)\,\Delta z$.}
	\label{fig:concio_eq_indefinite}
\end{figure}
